\chapter{Sprint 1}
\label{chap:sprint-1}

This chapter presents the realization of Sprint 1, the foundation every
later sprint builds on: turning an uploaded CV or TDR into structured,
extracted data end to end, proven with a working screen rather than just a
backend status update.

\section{Sprint Goal}

\begin{quote}
Upload a real CV/TDR and see it turned into structured, extracted data
end-to-end -- the foundation every later sprint builds on, proven with a
working (if rough) screen, not just a backend status update.
\end{quote}

Sprint 1 ran July 6--10, 2026, with 43 points committed -- deliberately
front-loaded, since it also had to prove the riskiest technical
integration of the whole project (multi-key LLM provider routing across
free-tier providers, since there is no paid API access) rather than only
low-complexity setup work.

\section{Sprint Backlog}

\begin{table}[h]
  \centering
  \caption{Sprint 1 backlog}
  \begin{tabular}{@{}p{1.3cm}p{7.7cm}p{1.3cm}p{2.5cm}@{}}
    \toprule
    ID & Story & Points & Status \\
    \midrule
    ADQ-1 & Provision Alembic migrations for the full Postgres schema & 3 & Done \\
    ADQ-2 & Multi-provider model routing: Groq + Gemini key-pool rotation with local Ollama fallback & 8 & Done \\
    ADQ-3 & Parse CV/TDR documents (text-PDF, docx) into raw text & 5 & Done \\
    ADQ-4 & Spike: sample real CVs/TDRs, confirm OCR need & 2 & Done \\
    ADQ-5 & Structured extraction for CV via Pydantic schema & 8 & Done \\
    ADQ-6 & Structured extraction for TDR + section localization & 8 & Done \\
    ADQ-7 & Bare-bones upload + extraction-preview page & 5 & Done \\
    ADQ-8 & Annotate gold extraction set (kickoff, ongoing) & 3 & Ongoing \\
    ADQ-9 & Request historical staffing decisions + TDR variety from company & 1 & Blocked (awaiting company response) \\
    \bottomrule
  \end{tabular}
\end{table}

Full acceptance criteria for every story are documented in
\texttt{docs/scrum/sprint-1-plan.md}, following the Definition of Ready and
Definition of Done in \texttt{docs/scrum/definition-of-ready-and-done.md}.

\section{Analysis}
% TODO: fill in once the sprint is fully wrapped up -- use-case-level
% analysis of the upload -> extraction -> preview flow.

\section{Design}
% TODO: Use Case Diagram, Class Diagram (ORM models / domain models),
% Sequence Diagram for the upload -> ingest -> extract -> persist flow.

\section{Implementation}
% TODO: important interfaces (ModelRouter, extract_cv/extract_tdr), code
% snippets only where they clarify a non-obvious decision, API endpoints
% table (POST /candidates/cv, POST /tdrs).

\section{Testing}
% TODO: summarize the unit test suite (key-pool rotation, ingest/extraction
% against real TDR samples and synthetic CVs, service-level tests), and
% note what is intentionally not yet covered (no live LLM integration test
% without real API keys).

\section{Sprint Review}
% TODO: describe the actual demo once given -- upload a real TDR and a
% synthetic CV through the running pages, show the extracted fields.

\section{Sprint Retrospective}
% TODO: what worked, what to change for Sprint 2 (e.g. was 43 points
% over-committed given the LLM routing risk called out in sprint planning).

\section{Conclusion}
% TODO: one short paragraph closing out the chapter once the above
% sections are filled in.
